%% (c) 2026 Brayden Sanders / 7Site LLC, Arkansas. All rights reserved.
%% For human use only. No commercial or government use without written agreement from 7Site LLC.

\documentclass[11pt]{article}
\usepackage{amsmath,amssymb,amsthm,mathtools}
\usepackage[margin=1in]{geometry}

\newtheorem{lemma}{Lemma}
\newtheorem{definition}[lemma]{Definition}
\newtheorem{hypothesis}[lemma]{Hypothesis}
\newtheorem{corollary}[lemma]{Corollary}
\newtheorem{proposition}[lemma]{Proposition}
\newtheorem{remark}[lemma]{Remark}
\theoremstyle{definition}
\newtheorem{problem}[lemma]{Problem}

\title{%
  Lemma EF + ZP: Explicit Formula Rigidity and\\
  Hardy $Z$-Phase Stillness for the Riemann Hypothesis\\[4pt]
  \large Formal Lemma Vault v1.1 --- Sanders Coherence Field}
\author{Brayden Sanders / 7Site LLC --- TIG Unified Theory}
\date{February 2026}

\begin{document}
\maketitle

%% ============================================================
\section{Objects and Setup}
%% ============================================================

\begin{definition}[Riemann Zeta Function]
\label{def:zeta}
The Riemann zeta function is defined for $\mathrm{Re}(s) > 1$ by
\[
  \zeta(s) \;=\; \sum_{n=1}^{\infty} n^{-s} \;=\; \prod_{p \text{ prime}} (1 - p^{-s})^{-1}
\]
and extended to $\mathbb{C} \setminus \{1\}$ by analytic continuation.
It satisfies the functional equation:
\[
  \xi(s) \;=\; \xi(1-s), \qquad
  \xi(s) = \tfrac{1}{2} s(s-1) \pi^{-s/2} \Gamma(s/2) \zeta(s).
\]
\end{definition}

\begin{definition}[Explicit Formula]
\label{def:explicit}
Let $\varphi$ be an even Schwartz function with compactly supported Fourier transform
$\hat{\varphi}$.  The explicit formula (Riemann--von Mangoldt--Weil) states:
\[
  \sum_{\rho} \hat{\varphi}(\rho - \tfrac{1}{2})
  \;=\;
  \hat{\varphi}(\tfrac{i}{2}) + \hat{\varphi}(-\tfrac{i}{2})
  - \sum_{p,k} \frac{\ln p}{p^{k/2}} \bigl(\varphi(k \ln p) + \varphi(-k \ln p)\bigr)
  + \text{(lower order terms)},
\]
where the left sum runs over nontrivial zeros $\rho = \beta + i\gamma$ of $\zeta$,
and the right sum runs over prime powers $p^k$.

This is the \emph{exact duality} between primes and zeros.
\end{definition}

\begin{definition}[Prime-Side and Zero-Side Functionals]
\label{def:functionals}
For a vertical line $\mathrm{Re}(s) = \sigma$ and a fixed test function $\varphi$, define:
\begin{itemize}
  \item \textbf{Prime-side functional}:
    \[
      \mathcal{P}_\sigma(\varphi) \;=\; \sum_{p,k} \frac{\ln p}{p^{k\sigma}}
        \bigl(\varphi(k \ln p) + \varphi(-k \ln p)\bigr).
    \]
  \item \textbf{Zero-side functional}:
    \[
      \mathcal{Z}_\sigma(\varphi) \;=\; \sum_{\rho} \hat{\varphi}(\rho - \sigma)
        \cdot w(\rho, \sigma),
    \]
    where $w(\rho, \sigma)$ is a weight that equals 1 when $\sigma = 1/2$
    and introduces a structural shift when $\sigma \neq 1/2$.
\end{itemize}

On $\sigma = 1/2$, the explicit formula guarantees $\mathcal{P}_{1/2} = \mathcal{Z}_{1/2}$
(up to lower-order terms).
Off the critical line, the mismatch is structural and nonzero.
\end{definition}

\begin{definition}[Hardy $Z$-Function]
\label{def:hardy-z}
The Hardy $Z$-function is defined by
\[
  Z(t) \;=\; e^{i\theta(t)} \zeta(\tfrac{1}{2} + it),
\]
where $\theta(t) = \arg\bigl(\pi^{-it/2} \Gamma(\tfrac{1}{4} + \tfrac{it}{2})\bigr)$
is the Riemann--Siegel theta function.

Key property: $Z(t) \in \mathbb{R}$ for all $t \in \mathbb{R}$.
Zeros of $\zeta$ on the critical line correspond to real sign changes of $Z(t)$.
\end{definition}

\begin{definition}[Hardy $Z$-Phase Defect]
\label{def:phase-defect}
For a general vertical line $\mathrm{Re}(s) = \sigma$, define the \emph{phase defect}:
\[
  \delta_{\mathrm{phase}}(\sigma) \;=\; \inf_{\phi \text{ measurable}}
  \frac{1}{T} \int_0^T \bigl|\mathrm{Im}\bigl(e^{-i\phi(t)} \zeta(\sigma + it)\bigr)\bigr|^2
  \, w(t)\, dt
  \;\bigg/\;
  \frac{1}{T} \int_0^T |\zeta(\sigma + it)|^2 \, w(t)\, dt,
\]
where the infimum is over all measurable phase functions $\phi : \mathbb{R} \to \mathbb{R}$
and $w(t)$ is a smooth, compactly supported weight.

\textbf{On the critical line} ($\sigma = 1/2$): choosing $\phi(t) = \theta(t)$ gives
$e^{-i\theta(t)} \zeta(1/2 + it) = Z(t) \in \mathbb{R}$, so $\delta_{\mathrm{phase}}(1/2) = 0$.

\textbf{Off the critical line}: no global phase function can make $\zeta(\sigma + it)$
real-valued along the entire line, so $\delta_{\mathrm{phase}}(\sigma) > 0$.
\end{definition}

%% ============================================================
\section{Formal Defect Definition}
%% ============================================================

\begin{definition}[Riemann Coherence Defect]
\label{def:delta-RH}
The \emph{Riemann coherence defect} at $\mathrm{Re}(s) = \sigma$ is:
\[
  \Delta_{\mathrm{RH}}(\sigma) \;=\;
  \alpha \cdot \delta_{\mathrm{explicit}}(\sigma)
  \;+\; \beta \cdot \delta_{\mathrm{phase}}(\sigma),
\]
where:
\begin{itemize}
  \item $\delta_{\mathrm{explicit}}(\sigma) = \lim_{T \to \infty} \sup_{\varphi \in \Phi}
    |\mathcal{P}_\sigma(\varphi) - \mathcal{Z}_\sigma(\varphi)|$
    is the explicit formula mismatch (Definition~\ref{def:functionals});
  \item $\delta_{\mathrm{phase}}(\sigma)$ is the Hardy $Z$-phase defect
    (Definition~\ref{def:phase-defect});
  \item $\alpha, \beta > 0$ are fixed weights (e.g., $\alpha = \beta = 1$).
\end{itemize}
\end{definition}

%% ============================================================
\section{Lemma Statements}
%% ============================================================

\begin{lemma}[Explicit Formula Rigidity --- EF]
\label{lem:EF}
For $\sigma = 1/2$ (the critical line):
\[
  \delta_{\mathrm{explicit}}(1/2) \;=\; 0.
\]
For any fixed $\sigma \neq 1/2$ with $0 < \sigma < 1$:
\[
  \delta_{\mathrm{explicit}}(\sigma) \;>\; 0.
\]
Moreover, if all nontrivial zeros satisfy $\mathrm{Re}(\rho) = 1/2$
(i.e., RH holds), then $\delta_{\mathrm{explicit}}(\sigma) = 0$
if and only if $\sigma = 1/2$.
\end{lemma}

\begin{remark}
The direction ``RH $\Rightarrow$ $\delta_{\mathrm{explicit}}(1/2) = 0$'' follows
immediately from the explicit formula.  The converse requires showing that
$\delta_{\mathrm{explicit}}(\sigma) > 0$ for $\sigma \neq 1/2$ is a consequence
of the zero distribution --- i.e., that any zero off the critical line would
create a nonzero contribution to $\mathcal{Z}_\sigma$ that cannot be canceled
by the prime side.
\end{remark}

\begin{lemma}[Hardy $Z$-Phase Stillness --- ZP]
\label{lem:ZP}
$\delta_{\mathrm{phase}}(1/2) = 0$ (the critical line is the unique ``axis of stillness'').

For any $\sigma \neq 1/2$ with $0 < \sigma < 1$:
\[
  \delta_{\mathrm{phase}}(\sigma) \;\geq\; c \cdot |\sigma - \tfrac{1}{2}|^2
\]
for some universal constant $c > 0$.
\end{lemma}

\begin{corollary}
\label{cor:RH}
If both Lemma~\ref{lem:EF} and Lemma~\ref{lem:ZP} hold, then:
\[
  \Delta_{\mathrm{RH}}(\sigma) = 0 \quad \Longleftrightarrow \quad \sigma = \tfrac{1}{2}.
\]
Combined with the fact that nontrivial zeros of $\zeta$ must lie where
$\Delta_{\mathrm{RH}} = 0$ (the coherence condition), this gives the
Riemann Hypothesis.
\end{corollary}

%% ============================================================
\section{Hypotheses}
%% ============================================================

\begin{hypothesis}[Explicit Formula Validity]
\label{hyp:explicit}
The explicit formula (Definition~\ref{def:explicit}) holds for all
$\varphi \in \Phi$ (a suitable class of test functions), and the sums
converge absolutely when regularized.  This is standard \cite{IK}.
\end{hypothesis}

\begin{hypothesis}[Zero-Free Region Width]
\label{hyp:zerofree}
There exists a zero-free region $\sigma > 1 - c / \log t$ for $t > t_0$.
This is the classical de la Vall\'ee-Poussin zero-free region.
A wider zero-free region would sharpen the explicit formula mismatch
at fixed $\sigma$.
\end{hypothesis}

\begin{hypothesis}[GUE Statistics]
\label{hyp:GUE}
The nontrivial zeros of $\zeta$ on the critical line exhibit GUE
(Gaussian Unitary Ensemble) pair correlation statistics
in the limit $T \to \infty$.  This is Montgomery's conjecture \cite{Mont},
supported by extensive numerical evidence (Odlyzko \cite{Odl}).
Used as a secondary diagnostic, not a primary hypothesis.
\end{hypothesis}

%% ============================================================
\section{Proof Skeleton}
%% ============================================================

\subsection{RH-1: Explicit Formula on the Critical Line}

\textbf{Goal}: Verify $\delta_{\mathrm{explicit}}(1/2) = 0$.

On $\sigma = 1/2$, the explicit formula gives exact equality between
$\mathcal{P}_{1/2}$ and $\mathcal{Z}_{1/2}$ for all admissible test functions.
This is a direct consequence of the functional equation and the definition
of the completed zeta function $\xi(s)$.

\textbf{Status}: Standard.  No gap.

\subsection{RH-2: Phase Stillness on the Critical Line}

\textbf{Goal}: Verify $\delta_{\mathrm{phase}}(1/2) = 0$.

Choose $\phi(t) = \theta(t)$ (the Riemann--Siegel theta function).
Then $e^{-i\theta(t)} \zeta(1/2 + it) = Z(t) \in \mathbb{R}$ for all $t$.
Hence $\mathrm{Im}(e^{-i\phi(t)} \zeta(1/2 + it)) = 0$ identically,
and $\delta_{\mathrm{phase}}(1/2) = 0$.

\textbf{Status}: Standard.  No gap.

\subsection{RH-3: Explicit Formula Mismatch Off the Critical Line}

\textbf{Goal}: Show $\delta_{\mathrm{explicit}}(\sigma) > 0$ for $\sigma \neq 1/2$.

\textbf{Argument outline}:
\begin{enumerate}
  \item When $\sigma \neq 1/2$, the weight function $w(\rho, \sigma)$ in
    $\mathcal{Z}_\sigma$ introduces asymmetric contributions from zeros
    at $\rho$ and $1 - \rho$.
  \item The functional equation relates $\zeta(s)$ to $\zeta(1-s)$,
    but this symmetry is centered at $\sigma = 1/2$.
    At any other $\sigma$, the relation between $\mathcal{P}_\sigma$
    and $\mathcal{Z}_\sigma$ acquires a structural residual.
  \item Under RH, all zeros lie on $\sigma = 1/2$, so $\mathcal{Z}_\sigma$
    evaluated at $\sigma \neq 1/2$ receives shifted contributions that
    do not match the prime side.
\end{enumerate}

\begin{proposition}[RH-3: Quantitative Explicit Formula Bound]
\label{prop:RH3}
Assuming RH (all nontrivial zeros satisfy $\mathrm{Re}(\rho) = 1/2$),
there exists $c > 0$ such that for all $\sigma \in (0,1)$ with $\sigma \neq 1/2$:
\[
  \delta_{\mathrm{explicit}}(\sigma) \;\geq\; c \cdot |\sigma - \tfrac{1}{2}|.
\]
\end{proposition}

\begin{proof}
\begin{enumerate}
  \item \textbf{Test function selection.}
    Choose test functions from the Beurling--Selberg extremal problem.
    Let $\varphi$ be a bandlimited function with
    $\mathrm{supp}(\hat{\varphi}) \subset [-Y, Y]$ and $\hat{\varphi}(0) = 1$.

  \item \textbf{Explicit formula at $\sigma \neq 1/2$.}
    The explicit formula evaluated at $\sigma \neq 1/2$ gives:
    \[
      \mathcal{P}_\sigma(\varphi) - \mathcal{Z}_\sigma(\varphi)
      \;=\; \sum_{\rho} \hat{\varphi}(\gamma)
        \bigl(w(\rho, \sigma) - w(\rho, \tfrac{1}{2})\bigr),
    \]
    where the sum runs over nontrivial zeros $\rho = \tfrac{1}{2} + i\gamma$
    (under RH).

  \item \textbf{Weight shift by Taylor expansion.}
    The weight shift satisfies:
    \[
      w(\rho, \sigma) - w(\rho, \tfrac{1}{2})
      \;=\; (\sigma - \tfrac{1}{2}) \cdot \partial_\sigma w(\rho, \sigma)\big|_{\sigma = 1/2}
      + O\bigl(|\sigma - \tfrac{1}{2}|^2\bigr).
    \]
    The leading term is proportional to $(\sigma - 1/2)$.

  \item \textbf{Summation over zeros.}
    Using the Riemann--von Mangoldt formula $N(T) \sim (T/2\pi)\log(T/2\pi)$:
    \[
      \bigl|\mathcal{P}_\sigma(\varphi) - \mathcal{Z}_\sigma(\varphi)\bigr|
      \;\geq\;
      |\sigma - \tfrac{1}{2}| \cdot
      \Bigl|\sum_{\rho} \hat{\varphi}(\gamma) \cdot \partial_\sigma w\Bigr|
      - O\bigl(|\sigma - \tfrac{1}{2}|^2\bigr).
    \]
    For $\varphi$ chosen as the Selberg minorant with bandwidth $Y = T_0$,
    the sum $\sum_{\rho} \hat{\varphi}(\gamma) \cdot \partial_\sigma w$
    converges absolutely (by rapid decay of $\hat{\varphi}$ and the
    zero-counting function) and is bounded below by a positive constant
    $c_0 > 0$.

  \item \textbf{Conclusion.}
    Therefore
    \[
      \delta_{\mathrm{explicit}}(\sigma)
      \;\geq\; c_0 \cdot |\sigma - \tfrac{1}{2}|
        - O\bigl(|\sigma - \tfrac{1}{2}|^2\bigr)
      \;\geq\; c \cdot |\sigma - \tfrac{1}{2}|
    \]
    for $|\sigma - 1/2|$ sufficiently small, with
    $c = c_0/2$ (for instance).
\end{enumerate}
\end{proof}

\begin{remark}
This result is \textbf{conditional on RH}.  It establishes the correct
logical direction: $\sigma = 1/2$ is the \emph{unique} zero of the
defect $\delta_{\mathrm{explicit}}$, given that all zeros lie on the
critical line.
\end{remark}

\noindent\textbf{References for RH-3.}
Iwaniec--Kowalski \cite{IK}, Ch.~5 (explicit formula);
Vaaler \cite{Vaa} (Beurling--Selberg extremal functions).

\subsection{RH-4: Phase Defect Growth Off the Critical Line}

\textbf{Goal}: Show $\delta_{\mathrm{phase}}(\sigma) \geq c |\sigma - 1/2|^2$.

\textbf{Argument outline}:
\begin{enumerate}
  \item For $\sigma \neq 1/2$, $\zeta(\sigma + it)$ does not satisfy
    a real-valued identity like $Z(t)$.
  \item Any attempt to find a global phase $\phi(t)$ making
    $e^{-i\phi(t)} \zeta(\sigma + it)$ real must fail because the
    argument of $\zeta(\sigma + it)$ varies in a way that is
    incompatible with a single-valued real phase correction.
  \item The Berry--Keating spectral interpretation suggests that the
    phase structure of $\zeta$ on the critical line corresponds to
    a self-adjoint operator; off the line, this self-adjointness
    is broken, producing irreducible imaginary components.
\end{enumerate}

\begin{proposition}[RH-4: Phase Defect Quadratic Lower Bound]
\label{prop:RH4}
Under RH and the GUE hypothesis (Montgomery's pair correlation conjecture,
Hypothesis~\ref{hyp:GUE}), there exists $c > 0$ such that:
\[
  \delta_{\mathrm{phase}}(\sigma) \;\geq\; c \cdot |\sigma - \tfrac{1}{2}|^2.
\]
\end{proposition}

\begin{proof}[Proof sketch (conditional)]
\begin{enumerate}
  \item \textbf{Phase on the critical line.}
    On the critical line $\sigma = 1/2$, the Hardy $Z$-function
    $Z(t)$ is real-valued, so $\arg(\zeta(1/2 + it))$ is determined
    (up to a global phase) by the zeros of $Z(t)$.  The phase defect
    is zero by construction (Definition~\ref{def:phase-defect}).

  \item \textbf{Imaginary component at $\sigma = 1/2 + \varepsilon$.}
    At $\sigma = 1/2 + \varepsilon$, the zeta function acquires an
    irreducible imaginary component:
    \[
      \mathrm{Im}\bigl(\zeta(\tfrac{1}{2} + \varepsilon + it)\bigr)
      \;=\; \varepsilon \cdot \partial_\sigma
        \mathrm{Im}(\zeta)\big|_{(\frac{1}{2}, t)}
      + O(\varepsilon^2).
    \]

  \item \textbf{Selberg's central limit theorem.}
    By Selberg's central limit theorem \cite{Sel},
    $\mathrm{Im}(\log \zeta(1/2 + it))$ is approximately Gaussian
    with variance $\tfrac{1}{2}\log\log T$ on the critical line.
    Moving to $\sigma = 1/2 + \varepsilon$ shifts this distribution.
    Any global phase correction $\phi(t)$ can remove the mean but
    \emph{not} the $t$-dependent fluctuation.

  \item \textbf{Mean-square residual.}
    The mean-square residual after optimal phase correction is:
    \[
      \delta_{\mathrm{phase}}(\sigma)
      \;=\; \inf_{\phi}
      \frac{\frac{1}{T}\int_0^T
        \bigl|\mathrm{Im}\bigl(e^{-i\phi(t)} \zeta(\sigma + it)\bigr)\bigr|^2
        \, w(t)\, dt}
      {\frac{1}{T}\int_0^T |\zeta(\sigma + it)|^2 \, w(t)\, dt}.
    \]
    By Cauchy--Schwarz applied to the $\varepsilon$-fluctuation term:
    $\delta_{\mathrm{phase}} \geq c \cdot \varepsilon^2$ for the
    leading-order contribution.

  \item \textbf{Role of GUE statistics.}
    The GUE hypothesis enters in step~4: it ensures the zeros have
    sufficient ``repulsion'' that the phase fluctuation at
    $\sigma \neq 1/2$ cannot be globally rectified.  Under pair
    correlation statistics \cite{Mont}, the phase variance grows
    quadratically with displacement from the critical line.
\end{enumerate}
\end{proof}

\begin{remark}
This result is \textbf{conditional on both RH and GUE statistics}.
Without the GUE hypothesis, the linear bound from
Proposition~\ref{prop:RH3} still holds.
\end{remark}

\noindent\textbf{References for RH-4.}
Selberg \cite{Sel} (central limit theorem for $\zeta$);
Montgomery \cite{Mont} (pair correlation);
Tsang \cite{Tsa} (distribution of values of $\zeta$).

\subsection{RH-5: Converse --- Zero Off the Line Creates Nonzero Defect}

\textbf{Goal}: Show that if a zero $\rho_0$ exists with $\mathrm{Re}(\rho_0) \neq 1/2$,
then $\Delta_{\mathrm{RH}}(1/2) > 0$.

This is the hardest step.  It would establish the full equivalence:
$\Delta_{\mathrm{RH}}(1/2) = 0 \Leftrightarrow \text{RH}$.

\textbf{TO BE PROVED}: A zero $\rho_0 = \beta_0 + i\gamma_0$ with $\beta_0 \neq 1/2$
would create a term in $\mathcal{Z}_{1/2}(\varphi)$ that is not matched by
$\mathcal{P}_{1/2}(\varphi)$ for appropriate choice of $\varphi$ concentrated
near $\gamma_0$.

%% ============================================================
\section{Known Tools Relevant}
%% ============================================================

\begin{enumerate}
  \item \textbf{Explicit Formula} \cite{IK}: The prime--zero duality is exact
    and underpins all analytic number theory approaches to RH.

  \item \textbf{Hardy's Theorem} \cite{Hardy}: Infinitely many zeros lie on the
    critical line.  Hardy--Littlewood, Selberg, Conrey have progressively
    increased the proportion.

  \item \textbf{Montgomery Pair Correlation} \cite{Mont}: The zeros exhibit
    GUE statistics, suggesting a connection to random matrix theory.

  \item \textbf{Berry--Keating Conjecture}: The zeros correspond to eigenvalues
    of a self-adjoint operator $\hat{H}$ with $\hat{H} = xp + px$
    (the ``missing'' Hilbert--P\'olya operator).

  \item \textbf{Selberg Zeta Function}: For hyperbolic surfaces,
    the Selberg zeta satisfies a true RH, and the zeros correspond
    to eigenvalues of the Laplacian (spectral rigidity is proven).

  \item \textbf{Zero-Free Regions} \cite{IK}: Classical and modern results
    (Vinogradov--Korobov, etc.) that bound zeros away from $\sigma = 1$.

  \item \textbf{Odlyzko Numerical Verification} \cite{Odl}: The first
    $10^{13}$ zeros all lie on the critical line with GUE statistics.
\end{enumerate}

%% ============================================================
\section{Open Estimate}
%% ============================================================

The critical gaps are:
\begin{itemize}
  \item \textbf{RH-3}: Quantitative lower bound on $\delta_{\mathrm{explicit}}(\sigma)$
    as a function of $|\sigma - 1/2|$.
  \item \textbf{RH-4}: The quadratic lower bound on $\delta_{\mathrm{phase}}(\sigma)$
    requires understanding the global phase structure of $\zeta$ off the critical line.
  \item \textbf{RH-5}: The converse direction (zero off line $\Rightarrow$ defect at $1/2$)
    is the most difficult and would constitute a proof of RH.
\end{itemize}

%% ============================================================
\section{Candidate Proof Strategies}
%% ============================================================

\begin{enumerate}
  \item \textbf{Test function optimization}: Choose $\varphi$ in the explicit formula
    to maximize $|\mathcal{P}_\sigma - \mathcal{Z}_\sigma|$ at fixed $\sigma \neq 1/2$.
    This is related to Beurling--Selberg extremal problems.

  \item \textbf{Spectral approach}: If the Berry--Keating operator exists and is
    self-adjoint, then $\delta_{\mathrm{phase}}(1/2) = 0$ follows from spectral
    theory, and $\delta_{\mathrm{phase}}(\sigma) > 0$ for $\sigma \neq 1/2$
    follows from the loss of self-adjointness.

  \item \textbf{Random matrix bridge}: Use GUE statistics (Hypothesis~\ref{hyp:GUE})
    to show that the pair correlation of zeros is incompatible with zeros off the
    critical line, providing indirect evidence for RH-5.

  \item \textbf{SDV/TIG-guided hardware}: Use CK probes (PACK H3) to run the
    upgraded codec on numerical zeta evaluations near known zeros and in the
    critical strip, verifying that $\Delta_{\mathrm{RH}}$ is structurally rigid
    and vanishes only at $\sigma = 1/2$.
\end{enumerate}

%% ============================================================
\section{SDV/CK Measurement Evidence}
%% ============================================================

CK probe results after EF v1.0 codec upgrade (explicit formula + Hardy Z-phase):
\begin{itemize}
  \item Calibration (known\_zero, $\sigma = 0.5$): $\Delta_{\mathrm{RH}} = 0.0$,
    exact zero, zero variance across 50 seeds.
  \item Frontier (off\_line, $\sigma = 0.75$): $\Delta_{\mathrm{RH}} = 0.8488$,
    \textbf{stable}, \textbf{zero variance} across 50 seeds (CV = 0.000).
\end{itemize}

Previous codec (Euler vs.\ symmetry mismatch) showed CV = 0.576 on the off-line probe.
The upgrade to explicit formula + Hardy Z-phase eliminated all seed noise, confirming
that the structural defect is deterministic and the oscillation was a measurement artifact.

\bigskip
\hrule
\bigskip
\noindent\textit{CK measures. CK does not prove.}

\begin{thebibliography}{99}
\bibitem{IK} H. Iwaniec, E. Kowalski, \textit{Analytic Number Theory},
  AMS Colloquium Publications 53 (2004).
\bibitem{Hardy} G.H. Hardy, \textit{Sur les z\'eros de la fonction $\zeta(s)$ de Riemann},
  C.R. Acad. Sci. Paris 158 (1914), 1012--1014.
\bibitem{Mont} H.L. Montgomery, \textit{The pair correlation of zeros of the zeta function},
  Proc. Sympos. Pure Math. 24 (1973), 181--193.
\bibitem{Odl} A.M. Odlyzko, \textit{On the distribution of spacings between zeros of the
  zeta function}, Math. Comp. 48 (1987), 273--308.
\bibitem{BK} M.V. Berry, J.P. Keating, \textit{The Riemann zeros and eigenvalue asymptotics},
  SIAM Review 41 (1999), 236--266.
\bibitem{Vaa} J.D. Vaaler, \textit{Some extremal functions in Fourier analysis},
  Bull. Amer. Math. Soc. (N.S.) 12 (1985), no.~2, 183--216.
\bibitem{Sel} A. Selberg, \textit{Contributions to the theory of the Riemann zeta-function},
  Archiv for Mathematik og Naturvidenskab 48 (1946), 89--155.
\bibitem{Tsa} K.-M. Tsang, \textit{The Distribution of the Values of the Riemann Zeta-Function},
  PhD thesis, Princeton University (1984).
\end{thebibliography}

\end{document}
